\begin{abstract}
We demonstrate that benchmark popularity correlates with larger generalization failures in machine learning models. Models trained on high-popularity datasets (CIFAR-10) exhibit generalization gaps 21 percentage points larger than those trained on low-popularity same-domain datasets (SVHN), despite identical architectures and training protocols. We find that popular benchmarks attract disproportionate research investment---high-use datasets receive 24.68$\times$ more optimization-focused papers ($p$=0.010)---consistent with a \emph{benchmark co-evolution} hypothesis where the ML ecosystem's optimization converges on dataset-specific patterns. However, testing the proposed mechanism, we find that modern architectures show \emph{lower} texture bias than legacy ones, ruling out texture exploitation as the causal pathway. Our results suggest that benchmark selection itself confounds model evaluation: progress measured on popular benchmarks may not reflect true generalization capability. We recommend complementing popular benchmarks with diverse, less-optimized alternatives to improve deployment reliability.
\end{abstract}
